%% sec_architecture.tex
%% ─────────────────────────────────────────────────────────────────────────────
%% Section 4 — Wavefunction-Token Cascade Architecture
%% ─────────────────────────────────────────────────────────────────────────────

\section{Wavefunction-Token Cascade Architecture}\label{sec:architecture}

\subsection{Pipeline overview}

The complete \ggttg matrix element computation is decomposed into a
\textbf{five-kernel cascade pipeline} spanning five consecutive \aie tiles in
a single row:
\begin{equation}
  \underbrace{\texttt{K1}}_{\text{WFGEN}}
  \xrightarrow{\text{8-WF token}}
  \underbrace{\texttt{K2a}}_{\text{FF\_DIAG\_A}}
  \xrightarrow{\text{8-WF token}}
  \underbrace{\texttt{K2b}}_{\text{FF\_DIAG\_B}}
  \xrightarrow{\text{8-WF token}}
  \underbrace{\texttt{K3}}_{\text{VVV\_AMP}}
  \xrightarrow{\text{5-WF token}}
  \underbrace{\texttt{K4}}_{\text{VVVV\_COLOR}}
  \label{eq:pipeline}
\end{equation}
Only the boundary kernels (\texttt{K1} and \texttt{K4}) have stream I/O; all
interior kernels are purely cascade-in/cascade-out.  This clean separation
means the cascade pipeline design is independent of the I/O topology.

\begin{figure*}[t]
\centering
\includegraphics[width=\textwidth]{Fig1}
\caption{\textbf{Fig.\,1} Single cascade pipeline (one row of 5 AIE tiles).
  Stage~1 (K1, WFGEN) receives the input PSP packet stream, generates 8 external
  wavefunctions, and transmits an 8-WF token over the cascade bus.
  Stages~2a and 2b (K2a/K2b) evaluate fermion-loop quark-gluon diagrams D2--D7
  and D8--D14 respectively, accumulating colour-flow amplitudes (JAMP) into the
  token.  Stage~3 (K3, VVV\_AMP) handles triple-gluon vertex diagrams D1, D12,
  and D15, reducing the token width from 8~WFs to 5~WFs.  Stage~4 (K4,
  VVVV\_COLOR) computes the four-gluon contact diagram D16 and performs the
  colour-matrix reduction, emitting the scalar ME$^2$ result as an output
  packet stream.  Cascade labels show token type (8-WF or 5-WF) and the
  cascade bus width is 384~bits at 1.25~GHz.}
\label{fig:single_pipe}
\end{figure*}

\subsection{Diagram-to-kernel assignment}\label{sec:diagram_assign}

Table~\ref{tab:diag_assign} shows which Feynman diagrams are evaluated in each
kernel and which HELAS functions are required.

\begin{table}[h]
\caption{Diagram-to-kernel assignment for \ggttg.  D16 has three
sub-diagrams (a, b, c) using four-gluon contact vertices.  D12 is evaluated
in K3 alongside the triple-gluon diagrams because its required vertex
\texttt{VVV1\_0} resides there.}
\label{tab:diag_assign}
\begin{tabular}{@{}llll@{}}
\toprule
Kernel & Diagrams & Key HELAS functions & JAMP contribution \\
\midrule
K1 & --- (wavefunction generation) &
  \texttt{ixxxxx}, \texttt{oxxxxx}, \texttt{vxxxxx}, \texttt{VVV1P0\_1}
  & initialise to 0 \\
K2a & D2, D3, D4, D5, D6, D7 &
  \texttt{FFV1\_0/1/2}
  & accumulate \\
K2b & D8, D9, D10, D11, D13, D14 &
  \texttt{FFV1\_0/1/2}
  & accumulate \\
K3 & D1, D12, D15 &
  \texttt{FFV1P0\_3}, \texttt{VVV1\_0}
  & accumulate \\
K4 & D16a, D16b, D16c &
  \texttt{VVVV1/3/4P0\_1}, \texttt{FFV1\_0}
  & accumulate + colour reduction \\
\bottomrule
\end{tabular}
\end{table}

\subsection{Token design}\label{sec:token}

All inter-kernel state is passed via a \emph{cascade token}---a structured
sequence of cascade beats encoding wavefunctions and the partial JAMP
accumulator.  Two token formats are used:

\paragraph{Extended 8-WF token (K1$\to$K2a$\to$K2b$\to$K3).}
Contains five external wavefunctions ($w_0$--$w_4$) plus three precomputed VVV
propagators ($w_{5,01} = \texttt{VVV1P0\_1}(w_0,w_1)$,
$w_{10} = \texttt{VVV1P0\_1}(w_1,w_4)$,
$w_{5,04} = \texttt{VVV1P0\_1}(w_0,w_4)$), and JAMP[6]:
\begin{equation}
  \underbrace{8 \times 2\ \text{beats}}_{\text{wavefunctions}}
  +\ \underbrace{2\ \text{beats}}_{\text{JAMP}[6]}
  = 18\ \text{cascade beats per helicity}.
  \label{eq:extended_token}
\end{equation}
Each beat transfers $4\times\texttt{cfloat} = 8$ single-precision complex
values (32\,bytes) via the 384-bit accumulator.  Wavefunctions use 6 of 8
lanes (lanes~0--5 carry the physics values; lanes~6--7 are zero-padded).

\paragraph{Standard 5-WF token (K3$\to$K4).}
Contains only the five external wavefunctions plus JAMP[6]:
\begin{equation}
  5 \times 2 + 2 = 12\ \text{cascade beats per helicity}.
  \label{eq:standard_token}
\end{equation}
K3 computes \texttt{FFV1P0\_3}$(w_3, w_2)$ locally and does not need to
forward it, reducing the token by 3 wavefunctions.

\subsection{Redundant propagator recomputation as a design choice}
\label{sec:redundant_recompute}

K2a and K2b each independently compute the same six propagators
($w_7, w_8, w_9, w_{5,20}, w_{5,30}, w_{11}$) from the external wavefunctions
carried in the token.  This is deliberate: passing the 6 propagators in the
token would require 12 additional cascade beats per helicity iteration,
increasing token size by 67\%.  If instead each propagator is recomputed
at a cost of ${\sim}200$ cycles, the cascade remains at 18 beats and tile
bandwidth is reduced by ${\sim}40\%$.  This computation-vs-communication
trade-off is central to keeping the architecture pipeline-bound rather than
cascade-bound.

\subsection{Deterministic cascade contract}\label{sec:contract}

To guarantee deadlock-free operation, all kernels in a pipeline must obey the
\textbf{cascade contract}:
\begin{enumerate}
  \item \textbf{Identical loop skeletons}: every kernel executes exactly
    \texttt{for (b in 0..BATCH) \{ for (h in 0..31) \{ ... \} \}}.  No early
    exits, no data-dependent iteration counts.
  \item \textbf{Unconditional cascade writes}: every iteration writes exactly
    one token to the cascade.  Conditional writes cause FIFO drift between
    adjacent tiles, eventually causing deadlock.
  \item \textbf{Statically matched counts}: the compiler-visible token
    write/read counts are equal across every producer/consumer pair at compile
    time.
\end{enumerate}
These rules are enforced by the token API implemented in
\texttt{kernels\_4k\_token.h}, which provides \texttt{write\_token\_8wf\_k1},
\texttt{read\_token\_8wf\_k1}, \texttt{write\_token\_5wf}, and
\texttt{read\_token\_5wf} as the only legal inter-kernel communication
primitives.

The 4-entry hardware FIFO on the cascade link absorbs minor cycle-to-cycle
timing skew (e.g., from branch mispredictions in the scalar unit) as long as
the loop-skeleton structure is preserved.

\subsection{Scalable 80-pipeline array architecture}\label{sec:array}

\paragraph{Group and row layout.}
The 400 \aie tiles are organised as $10\ \text{groups} \times 8\ \text{rows}$
$= 80$ pipelines.  Group~$g$ occupies columns $[5g, 5g+4]$; the five kernels
(\texttt{K1}--\texttt{K4}) occupy exactly those five columns.

\paragraph{Cascade direction compliance.}
Even rows cascade left-to-right; odd rows right-to-left.  Accordingly:
\begin{itemize}
  \item Even row $r$: \texttt{K1} at column $5g$, \texttt{K4} at $5g+4$.
  \item Odd row $r$: \texttt{K1} at column $5g+4$, \texttt{K4} at $5g$.
\end{itemize}
The HELAS kernel bodies are identical for all 80 pipelines; only the ADF
placement constraints differ.

\paragraph{Packet switching.}
Each group $g$ is served by one PLIO pair pinned to shim column $4g+6$
(giving columns 6, 10, 14, \ldots, 42---all valid shim sites on XCVC1902).  A
$\texttt{pktsplit}\langle 8\rangle$ node fans the incoming stream to all 8
\texttt{K1} tiles in the group (one per row); a $\texttt{pktmerge}\langle
8\rangle$ node collects all 8 \texttt{K4} outputs.  Packet ID bits
$[4:0]$ of the header encode the row index (0--7), routing each PSP to the
correct pipeline.  With 10 groups, the full array uses $10$ PLIO input pairs and
$10\times\texttt{pktsplit}\langle 8\rangle$.

\paragraph{Full-system three-layer stack.}
\begin{itemize}
  \item \textbf{PS (ARM A72)}: RAMBO phase-space generator produces PSPs;
    XRT runtime manages XCLBIN loading, buffer allocation, and synchronisation.
  \item \textbf{PL}: 10$\times$ HLS \texttt{MM2S\_pkt\_gen} kernels read PSP
    arrays from DDR and generate packet-headered AXI-Stream at II=1 (32-bit,
    500\,MHz).  10$\times$ HLS \texttt{S2MM\_pkt\_parser} collect
    \texttt{pktmerge} output and write ME$^2$ results to DDR.
  \item \textbf{AIE array}: 80 independent five-tile cascade pipelines
    consuming \texttt{input\_pktstream} in \texttt{K1} and emitting
    \texttt{output\_pktstream} from \texttt{K4}.
\end{itemize}

\begin{figure*}[t]
\centering
\includegraphics[width=\textwidth]{Fig2}
\caption{\textbf{Fig.\,2} Full SoC architecture of the VCK190 deployment.
  The Processing System (PS, ARM Cortex-A72, left) hosts the XRT host
  application and the RAMBO phase-space generator; DDR4 provides 32~GB of
  input/output memory.  The Programmable Logic (PL, orange, bottom-right)
  implements $10\times$ \texttt{MM2S\_pkt\_gen} and $10\times$
  \texttt{S2MM\_pkt\_parser} HLS kernels running at 500~MHz.  The AI Engine
  array (AIE, green, top-right) contains 80 five-tile cascade pipelines
  organised in 10 groups of 8 rows; each group is connected to one PLIO pair
  through \texttt{pktsplit}$\langle 8\rangle$ / \texttt{pktmerge}$\langle
  8\rangle$ nodes.  Bold arrows denote AXI-Stream PLIO links (32-bit, 500~MHz);
  the NOC provides AXI-MM access between PL and DDR.}
\label{fig:full_arch}
\end{figure*}

\subsection{Lazy evaluation: \texttt{FFV1P0\_3} moved to K3}
\label{sec:lazy_eval}

In the original partitioning, \texttt{FFV1P0\_3} was computed in \texttt{K1}
and passed in the token as wavefunction $w_6$.  PM profiling revealed that
\texttt{K1} with \texttt{FFV1P0\_3} included required ${\sim}17.8\KB$,
overflowing the 16\KB limit.  The function is needed only as an input to the
VVV amplitude diagrams D1, D12, and D15---all in \texttt{K3}.  Moving the
computation to \texttt{K3} (lazy evaluation) reduces \texttt{K1} to
${\sim}15.5\KB$ and simultaneously removes one wavefunction slot from the
token (from a 9-WF to an 8-WF token), saving 2 cascade beats per helicity
iteration.
